\scanpage{307}
\mspage{6}
s'agit de désigner \unsure{l'aspect en} ``objet simplicial'' ou ``complexe de chaînes'':
\[
 K_{\cdot}=ND_{\cdot}(K_*),\qquad K_*=DP_*(K_{\cdot})
\]
\[
 K_*=(n\mathrel{\unsure{\vdash|\mapsto}}K_{[n]}),\qquad
 K_{\cdot}=(K_n,d_n)_{n\geq0}
\]
On est ainsi amené à noter par une notation $K_{\cdot}$ tt complexe de chaînes,
par $K_n$ ses composantes (ce qui exclut l'utilisation de la notation $K_n$
pour la $n$.ième composante \inserted{de} $K_*$, et nous oblige à adopter
une notation telle que $K_{[n]}$!).

De même, le \unsure{th.ème|théorème} de $DP$ appliqué à $A^o$ nous donne des équivalences
\diagram{\begin{tikzcd}[column sep=large]
 A^*\arrow[r,shift left=1ex,"{ND^{\cdot}}","\approx"']&
 A^{\cdot}\arrow[l,shift left=1ex,"{DP^*}","\approx"']
\end{tikzcd}}
et des notations
\[
 K^{\cdot}=ND^{\cdot}(K^*),\qquad K^*=DP^*(K^{\cdot})
\]
\[
 K^{*}_{\deleted{\unsure{*}}}=(n\mathrel{\unsure{\vdash|\mapsto}}K^{[n]}),\qquad
 K^{\cdot}=(K^n,d^n)_{n\geq0}.
\]
Les foncteurs $ND^{\cdot}$ et $DP^*$ commutent aux foncteurs additifs.
En particulier, dans le cas d'une équivalence
\diagram{\begin{tikzcd}[column sep=large]
 A\arrow[r,leftrightarrow,"\approx","\vee"']&B^o
\end{tikzcd}}
déjà envisagé, étendant aux complexes de chaînes ou de cochaînes la convent.
de la montée ou de la descente des indices, on trouve que les foncteurs
$ND$ et $DP$ y commutent. Donc dans ce cas, la même lettre $K$ désigne
\margin{308}
